\section{Conclusão Parcial}

Até o encerramento da Sprint 1, a AGES IV tem representado uma mudança importante na minha atuação. Nas etapas anteriores, grande parte da minha preocupação estava em compreender tecnologias e concluir implementações. No ReadRace, além do trabalho técnico, passei a lidar mais diretamente com prioridades, dependências, acompanhamento de pessoas e comunicação com os stakeholders. Essa combinação ampliou minha percepção sobre o que significa contribuir para uma entrega coletiva.

Do ponto de vista técnico, participei de componentes compartilhados, funcionalidades mobile, regras de progresso de leitura e integração com o backend. A revisão da modelagem e das user stories também aproximou conceitos de requisitos, banco de dados e testes discutidos ao longo do curso. A regra que impede recompensar a mesma página mais de uma vez é um exemplo de como uma decisão de negócio precisa ser traduzida em dados, validações e comportamento consistente na interface.

O uso de agentes de IA como apoio à elaboração de tarefas, documentado no repositório de gestão, reforça uma reflexão iniciada na AGES III: compreender o negócio continua sendo essencial. Automatizar a escrita de uma tarefa não garante que a regra esteja correta, que o protótipo represente uma decisão validada ou que a dependência tenha sido resolvida. A equipe ainda precisa interpretar o problema, definir prioridades e avaliar o resultado. À medida que essas ferramentas passam a apoiar mais atividades, ganha relevância a capacidade de orientar o trabalho a partir do objetivo do produto.

Minha principal autocrítica está no acompanhamento da equipe. Fui solícito durante as aulas e procurei ajudar os colegas também fora delas, mas essas iniciativas não foram suficientes para evitar o acúmulo de pendências. Nós, AGES IV, precisávamos estar mais a par do andamento das tarefas e cobrar atualizações antes da proximidade do prazo. Assumir várias implementações ajudou a viabilizar a entrega, porém não resolve, por si só, o problema de gestão que levou à concentração do trabalho.

A retrospectiva permitiu transformar essa dificuldade em uma ação concreta: a divisão em squads com acompanhamento de um ou dois AGES IV. Minha expectativa é que essa organização favoreça a comunicação dos bloqueios e a distribuição das responsabilidades. Ao mesmo tempo, quero desenvolver um acompanhamento que ofereça apoio sem retirar a autonomia dos colegas. O resultado dessa mudança ainda precisa ser observado nas próximas sprints.

Outro aprendizado importante envolve a timidez. As apresentações para os stakeholders exigem que eu me exponha, organize o raciocínio e comunique resultados e limitações. Tenho participado ativamente ao lado dos demais AGES IV, e vejo nesse exercício uma oportunidade de ganhar confiança. Preparar os slides e conhecer o estado do projeto contribui para essa participação, mas é o contato com os clientes que me permite praticar essa habilidade.

Para a próxima etapa, pretendo acompanhar a equipe com mais constância, reduzir a concentração de tarefas no fim da sprint e continuar desenvolvendo minha comunicação. A avaliação do projeto como um todo dependerá das experiências e dos resultados das próximas entregas.
